%! Transformations of Functions — Unit 1, Lesson 1.4 — Student Packet Cover Sheet
\documentclass[10pt]{article}
\usepackage{precalculus-article}
\usepackage{precalculus-boxes}
\usepackage{ltablex}
\keepXColumns

\begin{document}

% Full-bleed plum banner
\begin{tikzpicture}[remember picture, overlay]
  \fill[plum] (current page.north west)
    rectangle ([yshift=-0.9in]current page.north east);
\end{tikzpicture}
\vspace{-0.8in}

\begin{center}
  {\color{white}\textbf{\LARGE Precalculus}}\\[4pt]
  {\color{lilacmid}\large\textbf{Unit 1: Functions and Change}}\\[6pt]
  {\color{white}\textbf{\large Lesson 1.4 \quad Transformations of Functions}}
\end{center}

\vspace{0.22in}
\namedateperiod
\vspace{0.18in}

\begin{learningtargetbox}
\small
\begin{enumerate}[leftmargin=1.5em, itemsep=5pt]
  \item I can tell a change made \textit{outside} a function from a change
        made \textit{inside} it, and predict whether the graph moves
        vertically or horizontally.
  \item I can build a table of values for $g(x) = af(x+h)+k$ from a parent
        function $f$, and describe each move in words --- shift, stretch,
        compress, or flip.
  \item I can find the domain and range of a transformed function from the
        domain and range of its parent.
\end{enumerate}
\end{learningtargetbox}

\vspace{0.14in}

\begin{tocbox}
\small
\renewcommand{\arraystretch}{1.5}
\begin{tabularx}{\linewidth}{@{} c l X r @{}}
\rowcolor{lilac}
\textbf{\#} & \textbf{Component} & \textbf{Description} & \textbf{Score} \\
\midrule
1 & Warm-Up        & Spiral review of prerequisite skills               & \blank{1.2cm} \\
2 & Guided Notes   & Shifts, stretches, and flips of a parent function  & \blank{1.2cm} \\
3 & Group Activity & The skate park ramp: redesigning one curve         & \blank{1.2cm} \\
4 & Exit Ticket    & Independent check on describing a transformation   & \blank{1.2cm} \\
5 & Homework       & Practice and extension                             & \blank{1.2cm} \\
\midrule
  & \multicolumn{2}{r}{\textbf{Total}} & \blank{1.2cm} \\
\end{tabularx}
\end{tocbox}

\vspace{0.14in}

\begin{remindbox}
\small
The last three lessons measured how a function changes. Today the
\textbf{whole graph} moves: one parent curve shifted, stretched, and flipped
into a family of new ones. One sentence carries the lesson --- a change
\textit{outside} the function is vertical and does what you expect; a change
\textit{inside} is horizontal and does the opposite.
\end{remindbox}

\end{document}
